\chapter{Statistiques}

\noindent\textit{Sondages, résultats scolaires, enquêtes de quartier : la statistique
permet d'organiser des données, de les résumer par quelques nombres (moyenne, médiane) et
de les représenter par des graphiques. Ce chapitre reprend le vocabulaire de base et les
outils indispensables pour traiter une série statistique.}
\vspace{8pt}

\connecteBanner

\section{Vocabulaire statistique}

\begin{definition}[Vocabulaire de base]
La \textbf{population} est l'ensemble sur lequel porte une étude ; un \textbf{individu} en
est un élément. Le \textbf{caractère} étudié peut être : \textbf{qualitatif} (valeurs non
numériques, ex. un moyen de transport), \textbf{quantitatif discret} (valeurs numériques
isolées, ex. un nombre d'enfants) ou \textbf{quantitatif continu} (valeurs regroupées en
classes, ex. une taille). L'\textbf{effectif} d'une valeur (ou modalité) est le nombre
d'individus qui la possèdent ; l'\textbf{effectif total} $N$ est le nombre total
d'individus. La \textbf{fréquence} d'une valeur est le quotient de son effectif par
l'effectif total (souvent exprimée en \%).
\end{definition}

\begin{exemple}
On interroge $200$ élèves d'un établissement de Yaoundé sur leur moyen de transport
habituel.
\begin{center}
\small
\begin{tabular}{lcccc}
\toprule
Moyen de transport & Bus & Vélo & Marche & Voiture \\
\midrule
Effectif & 80 & 50 & 40 & 30 \\
Fréquence & $40\,\%$ & $25\,\%$ & $20\,\%$ & $15\,\%$ \\
\bottomrule
\end{tabular}
\end{center}
Ici, le caractère « moyen de transport » est \textbf{qualitatif} ; l'effectif total est
$N=200$.
\end{exemple}

\begin{remarque}
La somme des effectifs vaut toujours l'effectif total $N$ ; la somme des fréquences vaut
toujours $1$ (ou $100\,\%$).
\end{remarque}

\section{Effectifs et fréquences cumulés}

\begin{definition}[Effectifs et fréquences cumulés croissants]
Pour un caractère quantitatif, l'\textbf{effectif cumulé croissant} (ECC) d'une valeur est
la somme des effectifs de cette valeur et de toutes les valeurs \textbf{inférieures}. La
\textbf{fréquence cumulée croissante} (FCC) se calcule de la même façon avec les
fréquences.
\end{definition}

\begin{exemple}
Une enquête est menée auprès de $25$ familles d'un quartier de Douala sur leur nombre
d'enfants.
\begin{center}
\small
\begin{tabular}{lccccc}
\toprule
Nombre d'enfants & 0 & 1 & 2 & 3 & 4 \\
\midrule
Effectif & 3 & 7 & 8 & 5 & 2 \\
Fréquence & $12\,\%$ & $28\,\%$ & $32\,\%$ & $20\,\%$ & $8\,\%$ \\
ECC & 3 & 10 & 18 & 23 & 25 \\
FCC & $12\,\%$ & $40\,\%$ & $72\,\%$ & $92\,\%$ & $100\,\%$ \\
\bottomrule
\end{tabular}
\end{center}
\end{exemple}

\begin{illus}
\begin{tikzpicture}
\begin{axis}[
  ybar,bar width=13pt,width=8cm,height=5.6cm,
  xlabel={\small Nombre d'enfants},ylabel={\small Effectif},
  xtick=data,ymin=0,ymax=9.5,axis lines=left,
  enlarge x limits=0.15,
  nodes near coords,nodes near coords style={font=\scriptsize},
  tick label style={font=\scriptsize}]
\addplot[fill=onbo,draw=onbo] coordinates {(0,3) (1,7) (2,8) (3,5) (4,2)};
\end{axis}
\end{tikzpicture}
\fig{Figure — Diagramme en bâtons : nombre d'enfants dans les 25 familles interrogées.}
\end{illus}

\begin{remarque}
Le dernier effectif cumulé croissant est toujours égal à l'effectif total $N$ ; la dernière
fréquence cumulée croissante vaut toujours $100\,\%$.
\end{remarque}

\section{Regroupement en classes}

\begin{definition}[Classes]
Lorsqu'un caractère quantitatif continu prend beaucoup de valeurs différentes, on les
regroupe en \textbf{classes} de la forme $[a\,;b[$. Le \textbf{centre de la classe},
$\frac{a+b}{2}$, sert de valeur représentative pour les calculs.
\end{definition}

\begin{exemple}
Tailles (en cm) de $40$ élèves d'un collège, regroupées en classes d'amplitude $10$ :
\begin{center}
\small
\begin{tabular}{lcccc}
\toprule
Classe (cm) & $[140\,;150[$ & $[150\,;160[$ & $[160\,;170[$ & $[170\,;180[$ \\
\midrule
Effectif & 6 & 14 & 15 & 5 \\
Fréquence & $15\,\%$ & $35\,\%$ & $37{,}5\,\%$ & $12{,}5\,\%$ \\
ECC & 6 & 20 & 35 & 40 \\
\bottomrule
\end{tabular}
\end{center}
\end{exemple}

\begin{illus}
\begin{tikzpicture}[scale=1]
\draw[onbmuted,->] (-0.3,0) -- (4.5,0);
\draw[onbmuted,->] (0,-0.2) -- (0,3.4);
\node[right,font=\scriptsize] at (4.5,0) {Taille (cm)};
\draw[fill=onbo!55,draw=onbo,line width=0.8pt] (0,0) rectangle (1,1.2);
\draw[fill=onbo!55,draw=onbo,line width=0.8pt] (1,0) rectangle (2,2.8);
\draw[fill=onbo!55,draw=onbo,line width=0.8pt] (2,0) rectangle (3,3.0);
\draw[fill=onbo!55,draw=onbo,line width=0.8pt] (3,0) rectangle (4,1.0);
\foreach \x/\lbl in {0/140,1/150,2/160,3/170,4/180}
  \node[below,font=\scriptsize] at (\x,0) {\lbl};
\node[above,font=\scriptsize] at (0.5,1.2) {6};
\node[above,font=\scriptsize] at (1.5,2.8) {14};
\node[above,font=\scriptsize] at (2.5,3.0) {15};
\node[above,font=\scriptsize] at (3.5,1.0) {5};
\end{tikzpicture}
\fig{Figure — Histogramme : tailles de 40 élèves, regroupées en classes de 10 cm.}
\end{illus}

\begin{remarque}
Quand toutes les classes ont la \textbf{même amplitude}, la hauteur de chaque rectangle de
l'histogramme est proportionnelle à l'effectif de la classe.
\end{remarque}

\section{Représentations graphiques}

\begin{propriete}[Choisir une représentation]
On utilise généralement : un \textbf{diagramme en bâtons} pour un caractère quantitatif
discret ; un \textbf{diagramme circulaire} (ou semi-circulaire) pour un caractère
qualitatif, en calculant l'angle de chaque secteur par $\text{angle}=\text{fréquence}\times360^{\circ}$ ;
un \textbf{histogramme} pour un caractère quantitatif continu regroupé en classes.
\end{propriete}

\begin{illus}
\begin{tikzpicture}[scale=0.9]
\draw[fill=onbo,draw=white,line width=1pt] (0,0) -- (90:1.8) arc (90:-54:1.8) -- cycle;
\draw[fill=onbgreen,draw=white,line width=1pt] (0,0) -- (-54:1.8) arc (-54:-144:1.8) -- cycle;
\draw[fill=onbblue,draw=white,line width=1pt] (0,0) -- (-144:1.8) arc (-144:-216:1.8) -- cycle;
\draw[fill=onbink,draw=white,line width=1pt] (0,0) -- (-216:1.8) arc (-216:-270:1.8) -- cycle;
\node[white,font=\scriptsize\bfseries] at (18:1.15) {40\%};
\node[white,font=\scriptsize\bfseries] at (-99:1.15) {25\%};
\node[white,font=\scriptsize\bfseries] at (-180:1.15) {20\%};
\node[white,font=\scriptsize\bfseries] at (-243:1.15) {15\%};
\foreach \x/\c/\t in {-2.7/onbo/Bus,-1.1/onbgreen/Vélo,0.5/onbblue/Marche,2.1/onbink/Voiture} {
  \draw[fill=\c,draw=none] (\x,-2.35) rectangle ++(0.22,0.22);
  \node[right,font=\scriptsize] at (\x+0.28,-2.24) {\t};
}
\end{tikzpicture}
\fig{Figure — Diagramme circulaire : moyen de transport des 200 élèves interrogés.}
\end{illus}

\begin{exemple}
Pour la catégorie « Bus » (fréquence $40\,\%$) : angle $=0{,}40\times360^{\circ}=144^{\circ}$.
Pour « Voiture » (fréquence $15\,\%$) : angle $=0{,}15\times360^{\circ}=54^{\circ}$.
\end{exemple}

\begin{remarque}
Dans un diagramme circulaire, la somme des angles de tous les secteurs vaut toujours
$360^{\circ}$ (comme la somme des fréquences vaut toujours $100\,\%$).
\end{remarque}

\section{Moyenne}

\begin{definition}[Moyenne d'une série statistique]
La \textbf{moyenne} d'une série de valeurs $x_1,x_2,\ldots,x_p$ d'effectifs
$n_1,n_2,\ldots,n_p$ et d'effectif total $N=n_1+n_2+\cdots+n_p$ est :
\[
\bar{x}=\frac{n_1x_1+n_2x_2+\cdots+n_px_p}{N}
\]
\end{definition}

\begin{exemple}
Reprenons le nombre d'enfants des $25$ familles (section précédente) :
\[
\bar{x}=\frac{3\times0+7\times1+8\times2+5\times3+2\times4}{25}=\frac{0+7+16+15+8}{25}=\frac{46}{25}=1{,}84.
\]
\end{exemple}

\begin{propriete}[Moyenne pondérée par des coefficients]
Lorsque chaque valeur $x_i$ est affectée d'un \textbf{coefficient} $c_i$ (et non d'un
effectif), la moyenne pondérée est $\bar{x}=\dfrac{c_1x_1+c_2x_2+\cdots+c_px_p}{c_1+c_2+\cdots+c_p}$.
\end{propriete}

\begin{exemple}
Notes obtenues à trois devoirs, avec leurs coefficients : Devoir 1 : $12$ (coeff. $2$) ;
Devoir 2 : $15$ (coeff. $1$) ; Devoir 3 : $9$ (coeff. $3$).
\[
\bar{x}=\frac{12\times2+15\times1+9\times3}{2+1+3}=\frac{24+15+27}{6}=\frac{66}{6}=11.
\]
\end{exemple}

\begin{remarque}
Une moyenne pondérée par des effectifs et une moyenne pondérée par des coefficients se
calculent exactement de la même façon : on somme les produits (valeur $\times$ poids), et on
divise par la somme des poids.
\end{remarque}

\section{Médiane et étendue}

\begin{definition}[Médiane et étendue]
La \textbf{médiane} d'une série ordonnée est une valeur qui partage la série en deux
groupes de même effectif : au moins la moitié des valeurs lui sont inférieures ou égales,
et au moins la moitié lui sont supérieures ou égales. L'\textbf{étendue} est la différence
entre la plus grande et la plus petite valeur de la série.
\end{definition}

\begin{propriete}[Calculer une médiane]
On \textbf{ordonne} d'abord la série. Si l'effectif total $N$ est \textbf{impair}, la
médiane est la valeur à la position $\dfrac{N+1}{2}$. Si $N$ est \textbf{pair}, la médiane
est la \textbf{moyenne} des deux valeurs aux positions $\dfrac{N}{2}$ et $\dfrac{N}{2}+1$.
\end{propriete}

\begin{illus}
\begin{tikzpicture}[scale=0.85]
\draw[onbline,line width=1pt] (0,0) -- (8,0);
\foreach \i/\v in {0/158,1/163,2/165,3/168,4/169,5/170,6/171,7/172,8/175} {
  \draw[onbline,line width=0.8pt] (\i,0.1)--(\i,-0.1);
  \node[below,font=\scriptsize] at (\i,-0.35) {\v};
}
\filldraw[onbo] (4,0) circle(4pt);
\node[above,onbo,font=\small\bfseries] at (4,0.35) {médiane};
\draw[onbmuted,<->] (0,-0.9) -- (3.85,-0.9);
\node[onbmuted,font=\scriptsize,below] at (2,-0.9) {4 valeurs};
\draw[onbmuted,<->] (4.15,-0.9) -- (8,-0.9);
\node[onbmuted,font=\scriptsize,below] at (6,-0.9) {4 valeurs};
\end{tikzpicture}
\fig{Figure — Série ordonnée de 9 tailles : la médiane (169) sépare deux groupes de 4 valeurs.}
\end{illus}

\begin{exemple}
Tailles (en cm) de $9$ joueurs, déjà ordonnées : $158$, $163$, $165$, $168$, $169$, $170$,
$171$, $172$, $175$.
$N=9$ est impair, la médiane est à la position $\frac{9+1}{2}=5$ : médiane $=169$ cm.
Étendue $=175-158=17$ cm.
\end{exemple}

\begin{remarque}
La médiane est \textbf{peu sensible} aux valeurs extrêmes, contrairement à la moyenne : elle
est souvent préférée pour résumer une série contenant des valeurs très éloignées des
autres.
\end{remarque}

% ═══════════════════════════════════════════════════════════════
\section{L'essentiel à retenir}

\begin{synthese}
\textbf{Vocabulaire.} Population, individu, caractère (qualitatif / quantitatif discret ou
continu), effectif, effectif total $N$, fréquence $=$ effectif $\div\, N$.\\[4pt]
\textbf{Cumulés.} ECC/FCC d'une valeur $=$ somme des effectifs/fréquences de cette valeur et
des valeurs inférieures ; dernier ECC $=N$, dernière FCC $=100\,\%$.\\[4pt]
\textbf{Représentations.} Bâtons (discret), circulaire (qualitatif, angle $=$ fréquence
$\times360^{\circ}$), histogramme (classes, amplitudes égales $\Rightarrow$ hauteur
$\propto$ effectif).\\[4pt]
\textbf{Moyenne.} $\bar{x}=\dfrac{\sum(\text{valeur}\times\text{poids})}{\sum\text{poids}}$
(poids $=$ effectifs ou coefficients).\\[4pt]
\textbf{Médiane.} Série ordonnée ; $N$ impair $\Rightarrow$ position $\frac{N+1}{2}$ ; $N$
pair $\Rightarrow$ moyenne des valeurs aux positions $\frac N2$ et $\frac N2+1$.\\[4pt]
\textbf{Piège fréquent.} Ne pas oublier d'\textbf{ordonner} la série avant de chercher la
médiane : la position dans le tableau brut n'a aucun rapport avec la médiane.
\end{synthese}

% ═══════════════════════════════════════════════════════════════
\section{Méthodes \& savoir-faire}

\methode{Construire un tableau d'effectifs et de fréquences à partir d'une liste brute}
Recenser toutes les valeurs différentes, compter combien de fois chacune apparaît
(effectif), puis diviser chaque effectif par l'effectif total pour obtenir la fréquence.
\begin{exemple}
Notes : $12$, $15$, $12$, $9$, $15$, $15$, $12$, $9$, $18$, $12$. Tableau : $9\to2$
($20\,\%$), $12\to4$ ($40\,\%$),
$15\to3$ ($30\,\%$), $18\to1$ ($10\,\%$) ; effectif total $=10$.
\end{exemple}

\methode{Calculer les effectifs et fréquences cumulés croissants}
Additionner progressivement les effectifs (ou les fréquences) de la plus petite valeur à
la plus grande.
\begin{exemple}
Pointures de $20$ personnes : $38\to2$, $39\to5$, $40\to7$, $41\to4$, $42\to2$. ECC :
$2$, $7$, $14$, $18$, $20$. FCC : $10\,\%$, $35\,\%$, $70\,\%$, $90\,\%$, $100\,\%$.
\end{exemple}

\methode{Calculer une moyenne pondérée par des effectifs}
Multiplier chaque valeur par son effectif, additionner ces produits, puis diviser par
l'effectif total.
\begin{exemple}
Valeurs $0$, $1$, $2$, $3$, $4$ d'effectifs $3$, $7$, $8$, $5$, $2$ (total $25$) :
$\bar{x}=\frac{0+7+16+15+8}{25}=\frac{46}{25}=1{,}84$.
\end{exemple}

\methode{Calculer une moyenne pondérée par des coefficients}
Multiplier chaque valeur par son coefficient, additionner ces produits, puis diviser par la
somme des coefficients (et non par le nombre de valeurs).
\begin{exemple}
Notes $13$ (coeff. $2$), $10$ (coeff. $1$), $16$ (coeff. $1$) :
$\bar{x}=\frac{13\times2+10\times1+16\times1}{2+1+1}=\frac{52}{4}=13$.
\end{exemple}

\methode{Déterminer une médiane}
Ordonner la série. Si l'effectif total est impair, prendre la valeur centrale. S'il est
pair, faire la moyenne des deux valeurs centrales.
\begin{exemple}
Série ordonnée de $8$ valeurs : $26$, $27$, $28$, $29$, $30$, $30$, $31$, $32$. Médiane
$=\frac{29+30}{2}=29{,}5$.
\end{exemple}

\methode{Calculer l'étendue d'une série}
Repérer la plus grande et la plus petite valeur de la série, puis calculer leur différence.
\begin{exemple}
Série $9$, $12$, $12$, $15$, $15$, $15$, $18$ : étendue $=18-9=9$.
\end{exemple}

\methode{Construire ou lire un diagramme circulaire}
Pour chaque catégorie, calculer l'angle du secteur par $\text{angle}=\text{fréquence}\times360^{\circ}$ ;
vérifier que la somme des angles vaut $360^{\circ}$.
\begin{exemple}
Fréquence $30\,\%$ : angle $=0{,}30\times360^{\circ}=108^{\circ}$. Fréquence $10\,\%$ :
angle $=0{,}10\times360^{\circ}=36^{\circ}$.
\end{exemple}

% ═══════════════════════════════════════════════════════════════
\section{Exercices d'application}
\begin{multicols}{2}

\rubrique{Tableaux d'effectifs et de fréquences}
\exo{}\diff{1} Pointures relevées chez $12$ personnes : $39$, $40$, $38$, $40$, $41$, $39$,
$40$, $38$, $40$, $39$, $41$, $40$. Construire le tableau des effectifs et des fréquences.

\exo{}\diff{1} Sport pratiqué par $50$ élèves : Football $20$, Basket $12$, Athlétisme $10$,
Natation $8$. Calculer les fréquences.

\exo{}\diff{2} Pointures de $20$ personnes : $38\to2$, $39\to5$, $40\to7$, $41\to4$,
$42\to2$. Calculer les effectifs et fréquences cumulés croissants.

\exo{}\diff{2} Nombre de livres lus pendant les vacances par $40$ élèves : $0\to4$,
$1\to9$, $2\to11$, $3\to8$, $4\to5$, $5\to3$. Calculer les effectifs et fréquences
cumulés croissants.

\rubrique{Moyennes}
\exo{}\diff{1} Nombre d'enfants de $25$ familles : $0\to3$, $1\to7$, $2\to8$, $3\to5$,
$4\to2$. Calculer la moyenne.

\exo{}\diff{2} Notes : $12$ (coeff. $2$), $15$ (coeff. $1$), $9$ (coeff. $3$). Calculer la
moyenne pondérée.

\exo{}\diff{2} Notes : $14$ (coeff. $3$), $8$ (coeff. $1$), $11$ (coeff. $2$). Calculer la
moyenne pondérée.

\exo{}\diff{2} Une entreprise camerounaise compte $5$ employés payés $150\,000$ FCFA, $3$
payés $200\,000$ FCFA et $2$ payés $300\,000$ FCFA. Calculer le salaire moyen.

\rubrique{Médiane et étendue}
\exo{}\diff{1} Tailles (en cm) de $9$ joueurs, déjà ordonnées :
$158$, $163$, $165$, $168$, $169$, $170$, $171$, $172$, $175$. Donner la médiane et
l'étendue.

\exo{}\diff{2} Pointures de $20$ personnes (exercice 3) : donner la médiane et
l'étendue.

\exo{}\diff{2} Nombre de livres lus par $40$ élèves (exercice 4) : donner la médiane et
l'étendue.

\exo{}\diff{2} Nombre d'enfants de $25$ familles (exercice 5) : donner la médiane et
l'étendue.

\rubrique{Représentations graphiques}
\exo{}\diff{1} Nombre de buts marqués lors de $20$ matchs : $0\to4$, $1\to6$, $2\to5$,
$3\to3$, $4\to2$. Calculer les fréquences (en vue d'un diagramme en bâtons).

\exo{}\diff{2} Mentions obtenues par $40$ élèves : Passable $10$, Assez bien $14$, Bien
$10$, Très bien $6$. Calculer les fréquences et les angles du diagramme circulaire.

\exo{}\diff{2} Poids (en kg) de $25$ nouveau-nés, par classes : $[2\,;2{,}5[\to3$,
$[2{,}5\,;3[\to9$, $[3\,;3{,}5[\to10$, $[3{,}5\,;4[\to3$. Calculer les fréquences (en vue
d'un histogramme).

\exo{}\diff{2} Nombre d'heures de sport par semaine, pour $50$ élèves : $0\to5$, $1\to10$,
$2\to15$, $3\to12$, $4\to6$, $5\to2$. Donner l'effectif total, le mode et l'étendue.

% ═══════════════════════════════════════════════════════════════
\end{multicols}

\section{Exercices d'entraînement}
\begin{multicols}{2}

\rubrique{Tableaux d'effectifs et de fréquences}
\exo{}\diff{1} Temps de trajet (en min) de $10$ employés : $15$, $20$, $15$, $25$, $20$,
$15$, $30$, $20$, $25$, $20$. Construire le tableau des effectifs et des fréquences.

\exo{}\diff{1} Sport pratiqué par $80$ élèves : Football $32$, Basket $20$, Handball $16$,
Autre $12$. Calculer les fréquences.

\exo{}\diff{2} Notes sur $10$ de $24$ élèves : $5\to2$, $6\to3$, $7\to6$, $8\to7$, $9\to4$,
$10\to2$. Calculer les effectifs et fréquences cumulés croissants.

\exo{}\diff{2} Nombre d'absences de $30$ élèves : $0\to12$, $1\to10$, $2\to5$, $3\to3$.
Calculer les effectifs et fréquences cumulés croissants.

\exo{}\diff{2} Animal préféré de $45$ élèves : Chat $18$, Chien $15$, Oiseau $7$, Poisson
$5$. Calculer les fréquences (arrondies au dixième).

\rubrique{Moyennes}
\exo{}\diff{2} Notes sur $10$ de $24$ élèves (exercice 3) : calculer la moyenne (arrondie
au centième).

\exo{}\diff{2} Notes : $16$ (coeff. $1$), $10$ (coeff. $3$), $13$ (coeff. $2$). Calculer la
moyenne pondérée.

\exo{}\diff{3} Une entreprise compte $8$ ouvriers payés $75\,000$ FCFA, $4$ techniciens
payés $125\,000$ FCFA et $3$ ingénieurs payés $250\,000$ FCFA. Calculer le salaire moyen
(arrondi au FCFA près).

\exo{}\diff{2} Nombre d'absences de $30$ élèves (exercice 4) : calculer la moyenne
(arrondie au centième).

\exo{}\diff{3} Tailles de $40$ élèves regroupées en classes de centres $145$, $155$, $165$,
$175$ et d'effectifs $6$, $14$, $15$, $5$. Calculer la moyenne (en utilisant les centres de
classes).

\rubrique{Médiane et étendue}
\exo{}\diff{2} Âges (en années) de $11$ personnes, ordonnés :
$11$, $12$, $12$, $13$, $13$, $14$, $14$, $15$, $15$, $16$, $17$. Donner la médiane et
l'étendue.

\exo{}\diff{2} Températures (en $^{\circ}$C) sur $10$ jours, ordonnées :
$19$, $20$, $21$, $22$, $23$, $24$, $25$, $26$, $27$, $28$. Donner la médiane et
l'étendue.

\exo{}\diff{2} Nombre d'absences de $30$ élèves (exercice 4) : donner la médiane et
l'étendue.

\exo{}\diff{2} Notes sur $10$ de $24$ élèves (exercice 3) : donner la médiane et
l'étendue.

\exo{}\diff{2} Nombre de buts marqués lors de $36$ matchs : $0\to6$, $1\to9$, $2\to10$,
$3\to7$, $4\to4$. Donner la médiane et l'étendue.

\rubrique{Représentations graphiques}
\exo{}\diff{1} Nombre de buts marqués lors de $36$ matchs (exercice précédent) : calculer
les fréquences.

\exo{}\diff{2} Sport préféré de $50$ personnes : Natation $15$, Football $20$, Tennis
$10$, Volley $5$. Calculer les fréquences et les angles du diagramme circulaire.

\exo{}\diff{2} Poids (en kg) de $60$ poissons pêchés, par classes : $[10\,;15[\to9$,
$[15\,;20[\to21$, $[20\,;25[\to18$, $[25\,;30[\to12$. Calculer les fréquences.

\exo{}\diff{2} Nombre de SMS envoyés par jour par $50$ élèves : $0\to4$, $1\to8$,
$2\to14$, $3\to12$, $4\to8$, $5\to4$. Donner le mode et l'étendue.

\exo{}\diff{3} Un sondage auprès de $300$ personnes sur leur fruit préféré donne : Mangue
$35\,\%$, Banane $25\,\%$, Ananas $20\,\%$, Orange $20\,\%$. Calculer l'effectif
correspondant à chaque fruit.

% ═══════════════════════════════════════════════════════════════
\end{multicols}

\section{Sujets type BEPC}

\sujetbac[Transport, notes et médiane]
\begin{enumerate}[label=\arabic*)]
\item Moyen de transport de $150$ personnes : Bus $60$, Taxi $45$, Moto $30$, Pied $15$.
Calculer les fréquences, puis les angles du diagramme circulaire correspondant.
\item Notes : $13$ (coeff. $2$), $10$ (coeff. $1$), $16$ (coeff. $1$). Calculer la moyenne
pondérée.
\item Âges (en années) de $7$ membres d'un club, ordonnés : $13$, $14$, $14$, $15$, $15$,
$16$, $17$. Donner la médiane et l'étendue.
\end{enumerate}

\sujetbac[Notes, absences et pointures]
\begin{enumerate}[label=\arabic*)]
\item Notes obtenues par $10$ élèves : $8$, $12$, $8$, $15$, $12$, $10$, $8$, $15$, $10$,
$12$. Construire le tableau des effectifs et des fréquences.
\item Nombre d'absences de $30$ élèves : $0\to12$, $1\to10$, $2\to5$, $3\to3$. Calculer
les effectifs et fréquences cumulés croissants.
\item Nombre d'enfants de $25$ familles : $0\to3$, $1\to7$, $2\to8$, $3\to5$, $4\to2$.
Calculer la moyenne.
\item Pointures de $20$ personnes : $38\to2$, $39\to5$, $40\to7$, $41\to4$, $42\to2$.
Donner la médiane.
\end{enumerate}

\sujetbac[Problème concret : salaires et poids de naissance]
\begin{enumerate}[label=\arabic*)]
\item Une entreprise compte $5$ employés payés $150\,000$ FCFA, $3$ payés $200\,000$ FCFA
et $2$ payés $300\,000$ FCFA. Calculer le salaire moyen.
\item Poids (en kg) de $25$ nouveau-nés, par classes : $[2\,;2{,}5[\to3$,
$[2{,}5\,;3[\to9$, $[3\,;3{,}5[\to10$, $[3{,}5\,;4[\to3$. Calculer les fréquences.
\item Températures (en $^{\circ}$C) sur $8$ jours, ordonnées : $26$, $27$, $28$, $29$,
$30$, $30$, $31$, $32$. Donner la médiane et l'étendue.
\item Mentions obtenues par $40$ élèves : Passable $10$, Assez bien $14$, Bien $10$, Très
bien $6$. Calculer les angles du diagramme circulaire.
\end{enumerate}

% ═══════════════════════════════════════════════════════════════
\section{Corrigés}
\begin{multicols}{2}

\corrige{de l'exercice 1}
$38\to2$ ($16{,}7\,\%$), $39\to3$ ($25\,\%$), $40\to5$ ($41{,}7\,\%$), $41\to2$
($16{,}7\,\%$) ; total $12$.

\corrige{de l'exercice 2}
Fréquences : Football $40\,\%$, Basket $24\,\%$, Athlétisme $20\,\%$, Natation $16\,\%$.

\corrige{de l'exercice 3}
ECC : $2$, $7$, $14$, $18$, $20$. FCC : $10\,\%$, $35\,\%$, $70\,\%$, $90\,\%$, $100\,\%$.

\corrige{de l'exercice 4}
ECC : $4$, $13$, $24$, $32$, $37$, $40$. FCC : $10\,\%$, $32{,}5\,\%$, $60\,\%$, $80\,\%$,
$92{,}5\,\%$, $100\,\%$.

\corrige{de l'exercice 5}
$\bar{x}=\dfrac{0+7+16+15+8}{25}=\dfrac{46}{25}=1{,}84$.

\corrige{de l'exercice 6}
$\bar{x}=\dfrac{12\times2+15\times1+9\times3}{6}=\dfrac{66}{6}=11$.

\corrige{de l'exercice 7}
$\bar{x}=\dfrac{14\times3+8\times1+11\times2}{6}=\dfrac{72}{6}=12$.

\corrige{de l'exercice 8}
$\bar{x}=\dfrac{5\times150\,000+3\times200\,000+2\times300\,000}{10}$\\
$=\dfrac{1\,950\,000}{10}=195\,000$ FCFA.

\corrige{de l'exercice 9}
Médiane $=169$ cm (position $5$ sur $9$). Étendue $=175-158=17$ cm.

\corrige{de l'exercice 10}
$N=20$ pair : médiane $=\dfrac{40+40}{2}=40$ (positions $10$ et $11$, toutes deux dans la
valeur $40$). Étendue $=42-38=4$.

\corrige{de l'exercice 11}
$N=40$ pair : les positions $20$ et $21$ correspondent toutes deux à la valeur $2$.
Médiane $=2$. Étendue $=5-0=5$.

\corrige{de l'exercice 12}
$N=25$ impair : position $13$, qui correspond à la valeur $2$ (ECC : $3$, $10$, $18$).
Médiane
$=2$. Étendue $=4-0=4$.

\corrige{de l'exercice 13}
Fréquences : $0\to20\,\%$, $1\to30\,\%$, $2\to25\,\%$, $3\to15\,\%$, $4\to10\,\%$.

\corrige{de l'exercice 14}
Fréquences : $25\,\%$, $35\,\%$, $25\,\%$, $15\,\%$. Angles : $90^{\circ}$, $126^{\circ}$,
$90^{\circ}$, $54^{\circ}$.

\corrige{de l'exercice 15}
Fréquences : $12\,\%$, $36\,\%$, $40\,\%$, $12\,\%$.

\corrige{de l'exercice 16}
Effectif total $=50$. Mode $=2$ heures (effectif $15$). Étendue $=5-0=5$.

\corrige{du sujet 1}
1) Fréquences : $40\,\%$, $30\,\%$, $20\,\%$, $10\,\%$. Angles : $144^{\circ}$, $108^{\circ}$,
$72^{\circ}$, $36^{\circ}$.\\
2) $\bar{x}=\dfrac{13\times2+10\times1+16\times1}{4}=\dfrac{52}{4}=13$.\\
3) $N=7$ impair, position $4$ : médiane $=15$. Étendue $=17-13=4$.

\corrige{du sujet 2}
1) $8\to3$ ($30\,\%$), $10\to2$ ($20\,\%$), $12\to3$ ($30\,\%$), $15\to2$ ($20\,\%$) ;
total $10$.\\
2) ECC : $12$, $22$, $27$, $30$. FCC : $40\,\%$, $73{,}3\,\%$, $90\,\%$, $100\,\%$.\\
3) $\bar{x}=\dfrac{46}{25}=1{,}84$.\\
4) Médiane $=40$ (positions $10$ et $11$).

\corrige{du sujet 3}
1) $\bar{x}=\dfrac{1\,950\,000}{10}=195\,000$ FCFA.\\
2) Fréquences : $12\,\%$, $36\,\%$, $40\,\%$, $12\,\%$.\\
3) $N=8$ pair : médiane $=\dfrac{29+30}{2}=29{,}5$. Étendue $=32-26=6$.\\
4) Angles : $90^{\circ}$, $126^{\circ}$, $90^{\circ}$, $54^{\circ}$.

\end{multicols}
\exercicesBanner
